\section{Arena v7: Temporal Multimodal Action and Hierarchical Sensing}
\label{sec:arena-v7-temporal-multimodal}

Arena v7 extends the natural-pixel grounding of Arena v6 into ordered visual
and acoustic time.  The evaluated cohort contains two official NASA videos and
one official USGS seismic recording.  Every source and derived artifact is
SHA--256 committed.  Perceptual proposals are committed before official source
claims are revealed, and unsupported exact identities are quarantined rather
than entering accepted knowledge.

\subsection{Ordered visual prediction}

Video is represented by six ordered frames.  Four frames are supplied for
proposal and two remain withheld.  The replaceable vision substrate proposes
persistent entities, observed change, a temporal process and a qualitative
next-frame prediction.  HexCore remains the evaluator and authority.  For a
frame order $\pi$, chronology consistency is measured by
\begin{equation}
 C(\pi)=\frac{1}{T-1}\sum_{t=1}^{T-1}
 \operatorname{MAE}\!\left(I_{\pi(t+1)},I_{\pi(t)}\right).
\end{equation}
A deliberately broken chronology must yield a greater cost than the source
order.  If the ordering is inconsistent, the system abstains from forecasting.
Arena v7 detected both broken chronologies and preserved 100\% mean and
weakest-case temporal semantic accuracy while observing four of six frames.

Official metadata was used only after proposal commitment.  The system also
derived and verified the temporal densities
\begin{equation}
 \frac{500\ \mathrm{frames}}{25\ \mathrm{hours}}=20\
 \ \mathrm{frames\,hour}^{-1},
 \qquad
 \frac{200\ \mathrm{steps}}{5\ \mathrm{hours}}=40\
 \ \mathrm{steps\,hour}^{-1}.
\end{equation}

\subsection{Audio representation and metacognitive measurement}

The audio path is explicitly not described as native listening.  The USGS
recording is represented by a spectrogram for semantic proposal and by
five-second root-mean-square (RMS) waveform windows for quantitative
measurement.  For samples $x_1,\ldots,x_n$, window energy is
\begin{equation}
 \operatorname{RMS}(x)=\sqrt{\frac{1}{n}\sum_{i=1}^{n}x_i^2}.
\end{equation}

The first governed run allocated all four detailed inspections inside the
spectrogram-proposed regions.  Although spectral interpretation was correct,
it recovered only 25\% of the four strongest events; CAU therefore rejected
the challenger.  The retained revision introduced hierarchical sensing:
two inspections exploit the semantic hypothesis and two explore the strongest
unexplained windows after a cheap full-stream energy scan.  With $N=39$,
coarse-scan cost $c_s=0.05$, detailed cost $c_d=1$, and $k=4$ detailed
inspections,
\begin{equation}
 C_{\mathrm{active}}=Nc_s+kc_d=5.95,
 \qquad C_{\mathrm{exhaustive}}=Nc_d=39,
\end{equation}
giving an 84.74\% measurement-cost reduction.  Strongest-event coverage rose
from 25\% to the 75\% promotion threshold.  This is a direct correction of
confirmation bias: the system reserves measurement capacity for observations
that its current interpretation does not explain.

\subsection{Governed outcome}

\begin{table}[ht]
\centering
\small
\begin{tabular}{@{}lr@{}}
\hline
Measure & Result \\
\hline
Public temporal sources / authorities & $3/3$ \\
Video temporal semantic accuracy & $100\%$ \\
Weakest video accuracy & $100\%$ \\
Broken chronologies detected & $2/2$ \\
Derived temporal quantities correct & $2/2$ \\
Seismic spectral semantic accuracy & $100\%$ \\
Strongest-event coverage & $75\%$ \\
Audio measurement-cost reduction & $84.74\%$ \\
Video observation reduction & $33.33\%$ \\
Governed temporal project success & $100\%$ \\
Unsafe or quarantined identity acceptances & $0$ \\
Restart relearning & $0$ \\
\hline
\end{tabular}
\caption{Sealed Arena v7 outcome.}
\end{table}

CAU promoted
\path{procedure_temporal_multimodal_action_v7_2a4710a57e8d}.
The temporal cohort, proposal commitments, selected procedure and authority
decision survive reconstruction with zero source relearning.

\subsection{Claim boundary}

This result establishes bounded temporal grounding, chronology criticism and
budgeted evidence allocation over three public natural-media sources.  Video is
sampled into frames and audio is represented through waveform energy and a
spectrogram.  Expected semantics, perturbations and evaluation remain
development controlled.  The result is not physical actuation, unrestricted
video/audio understanding, independent hidden certification or AGI.
